% --- Άσκηση 13319 (13319.tex) ---
\Askhsh{13319}{2}
\begin{ylh}{2}
    \item 2.1. Οι Πράξεις και οι Ιδιότητές τους
    \item 3.3. Εξισώσεις 2ου Βαθμού
    \item 5.2. Αριθμητική πρόοδος
    \item 5.3. Γεωμετρική πρόοδος
\end{ylh}
Δίνονται οι αριθμοί $1 - x\ ,\ \ \dfrac{x}{2},\ \ 2x\ –\ 1$, $x \in R$.
\begin{alist}
\item Να αποδείξετε ότι οι παραπάνω αριθμοί, με αυτή τη σειρά, είναι πάντοτε διαδοχικοί όροι μιας αριθμητικής προόδου.
\monades{13}
\item Να βρείτε την τιμή του $x$, αν γνωρίζουμε ότι η διαφορά ω αυτής της προόδου είναι 5.
\monades{12}
\end{alist}
